\chapter{应用实例}
\intro{本章展示数学物理方程在不同领域的应用。}

\section{热传导问题}

\begin{theorem}[有界杆的热传导]
\[
\begin{cases}
u_t=ku_{xx},&0<x<L\\
u(0,t)=u(L,t)=0\\
u(x,0)=\phi(x)
\end{cases}
\]
解为：
\[
u(x,t)=\sum_{n=1}^{\infty}b_n\sin\frac{n\pi x}{L}e^{-n^2\pi^2kt/L^2}
\]
\end{theorem}

\section{弦振动}

\begin{theorem}[固定端点的弦振动]
\[
u(x,t)=\sum_{n=1}^{\infty}\left(a_n\cos\frac{n\pi at}{L}+b_n\sin\frac{n\pi at}{L}\right)\sin\frac{n\pi x}{L}
\]
\end{theorem}

\section{静电场}

\begin{theorem}[球内 Dirichlet 问题]
\[
\begin{cases}
\lapl u=0,&|\vect{x}|<R\\
u|_{|\vect{x}|=R}=f(\theta,\phi)
\end{cases}
\]
解用球谐函数展开。
\end{theorem}

\begin{remark}
同一方程在不同边界条件下描述完全不同的物理现象，这体现了数学物理方程的统一性和普适性。
\end{remark}
